\documentclass[11pt]{letter}
\usepackage[utf8]{inputenc}
\usepackage[T1]{fontenc}
\usepackage{lmodern}
\usepackage[margin=1.1in]{geometry}
\usepackage{amsmath,amssymb}
\usepackage{microtype}
\usepackage{hyperref}
\hypersetup{hidelinks}

\address{David H.\ Silver\\
Independent researcher\\
Poughkeepsie, NY 12603, USA\\
\href{mailto:silverdavi@gmail.com}{silverdavi@gmail.com}\\
ORCID: \href{https://orcid.org/0000-0002-3071-304X}{0000-0002-3071-304X}}
\signature{David H.\ Silver}
\begin{document}
\begin{letter}{Editor-in-Chief\\ \emph{Annals of Applied Statistics}\\ Institute of Mathematical Statistics}

\opening{Dear Editor,}

I am submitting the enclosed manuscript, \emph{``Bounds and Estimators for the Civilian-to-Combatant Death Ratio from Aggregation of Conflicting Sources,''} for consideration in the \emph{Annals of Applied Statistics}.

The paper addresses a long-standing methodological problem in conflict statistics: how to draw quantitative inferences about civilian-to-combatant death ratios when the available sources have known, opposing political incentives. Each belligerent has a structural reason to inflate enemy combatant deaths and deflate civilian deaths, and intergovernmental and advocacy bodies have symmetric incentives in the other direction. Recent applied-statistics work in this area---multiple systems estimation (Lum, Price and Banks, 2013; Khatib et al., 2025), change-point modelling (Fagan et al., 2020), Bayesian aggregation under reporting bias (Radford et al., 2023), and demographic Bayesian inference for Gaza (G\'{o}mez-Ugarte et al., 2025)---has tackled the problem one channel at a time. The present manuscript instead treats the parameter of interest as \emph{partially identified} (Manski, 2003; Molinari, 2020; Kline and Tamer, 2023) and asks what set of values is jointly consistent with all sources, before any single channel is trusted.

The paper makes four methodological contributions:
\begin{enumerate}
\item A sharp identified set for the combatant share $q=D_M/D$ under a one-parameter exposure model that bounds the differential death rate of civilian adult men relative to women and children (Theorem 3.4); intersected with the mechanical manpower bound $D_M\le M$, this is a closed-form interval (Corollary 3.5).
\item A precision-weighted aggregation rule for multiple demographic sources, with delta-method standard error and a contamination-aware bias bound (Proposition 2.1).
\item A \emph{contradiction radius} $\rho$ that scores any disputed claim as the smallest $\ell_\infty$ standardised perturbation of independent inputs needed to make the claim feasible (Definition 4.2, Theorem 4.3).
\item A bias-robust credible interval under Huber $\varepsilon$-contamination of each source, with a sharper quantile-displacement bound under a density lower bound (Proposition 5.1)---building on Berger and Berliner (1986) and the recent revival of $\varepsilon$-contamination methods in Gagnon (2023) and Di Noia, Ruggeri and Mira (2025).
\end{enumerate}

The framework is applied in detail to the Gaza war 2023--2026, aggregating six independent source channels (PCBS, OHCHR, Gaza MoH, OCHA, IISS/RUSI, and Lancet field surveys). The public Israel Defense Forces claim of $17$--$25$k combatants killed has $\rho\gtrsim 20$ standard errors, while a spatial Bayesian posterior on the same independently-measured inputs gives $q=2.5\%$ with 95\% credible interval $[1.6\%,3.8\%]$ and an implied civilian-to-combatant ratio of $43{:}1$ \mbox{$[30,64]$}. Crucially, the test is symmetric: it rejects the polar narrative $q=0$ for the same reason---it is also outside the feasible set given the observed overrepresentation of adult males.

The paper fits AOAS's tradition of methodologically rigorous applied work. The closest precedents in your journal are the Bayesian hidden Markov model for ceasefire dynamics (vol.~18, 2024) and the Syrian conflict entity-resolution paper currently in your pipeline. The manuscript is written so that the methodological core (Sections 2--5) can be read independently of the Gaza application (Section 6) and the 81-conflict survey (Section 7, in the supplementary material). All proofs are stated as theorems; full proofs of the regularity-conditions extensions are in the supplement.

The manuscript runs approximately $9{,}500$ words plus a $20$-page supplement. Code, simulator, and Gaza inputs will be deposited in a public repository on acceptance. This work has not been considered or submitted elsewhere, and I declare no conflicts of interest.

\textbf{About the author.} I am an independent researcher; my peer-reviewed publications appear in \emph{Nature}, \emph{PNAS}, \emph{IEEE Trans.\ Pattern Analysis and Machine Intelligence}, \emph{Human Reproduction}, and \emph{PLOS ONE}. I held a Microsoft Research PhD Fellowship in Biomathematics, Bioinformatics, and Computational Biology jointly between the Technion and Microsoft Research Cambridge (UK), and I am the author of the peer-reviewed monograph \emph{Beyond Popular Science} (Open Book Publishers, Cambridge, 2026; DOI 10.11647/OBP.0526). The present manuscript is conducted independently and is unaffiliated with any of my industry roles.

Possible Associate Editors with the right combination of partial-identification and applied expertise include Charles Manski (Northwestern), Aaron Schein (Chicago), Tirthankar Dasgupta (Rutgers), and Susan Murphy (Harvard). Possible reviewers include Patrick Ball (HRDAG), Megan Price (HRDAG), Michael Spagat (Royal Holloway), Diego Alburez-Gutierrez (MPIDR), and Francesca Molinari (Cornell).

\closing{Sincerely,}

\end{letter}
\end{document}
